\documentclass[10pt]{article}
\usepackage[margin=0.72in]{geometry}
\usepackage{booktabs}
\usepackage{graphicx}
\usepackage{float}
\usepackage[hidelinks]{hyperref}
\hypersetup{pdftitle={Spine RQ3: Realized Work Explains Dynamic-Update Latency}}
\title{\textbf{Spine RQ3: Realized Work and Dynamic-Update Latency}\\
\large Correctness-Gated Cycle-Level Evidence}
\author{Chuxiao Han}
\date{July 2026}
\IfFileExists{data/rq3/representative_table.tex}{
  \newcommand{\rqdatadir}{data/rq3}
  \newcommand{\rqfigdir}{../figures}
}{
  \newcommand{\rqdatadir}{docs/paper/data/rq3}
  \newcommand{\rqfigdir}{docs/figures}
}
\begin{document}
\maketitle

\begin{abstract}
This report tests whether Spine latency follows work realized by the actual
execution rather than graph capacity or a fitted edge-count proxy. It covers
54 dual-oracle-admitted executions and all five requested mechanisms:
zero-net, shallow insertion, deep carry, residual PageRank correction, and
SSSP/CC nonmonotonic fallback. Every plotted latency ledger closes exactly to
end-to-end cycles.
\end{abstract}

\section{Representative end-to-end breakdown}
\begin{figure}[H]
  \centering
  \includegraphics[width=0.98\linewidth]{\rqfigdir/rq3_latency_breakdown.pdf}
  \caption{Exclusive, overlap-aware critical-path composition. Absolute
  device cycles are printed above each normalized bar. Integrated
  resolve/app is used only when the runner does not expose split timestamps.}
\end{figure}

\begin{table}[H]
\centering
\small
\input{\rqdatadir/representative_table.tex}
\caption{Absolute timing for the selected representative of each class.
Selection is deterministic: the maximum-cycle eligible execution.}
\end{table}

\clearpage
\section{Realized-work validation}
\begin{figure}[H]
  \centering
  \includegraphics[width=0.98\linewidth]{\rqfigdir/rq3_work_correlations.pdf}
  \caption{Direct execution counters versus mechanism-local cycles. Blue
  squares are calibration cases and red circles are holdout traces. Fits use
  all admitted rows with positive supported work.}
\end{figure}

\begin{table}[H]
\centering
\small
\input{\rqdatadir/regression_table.tex}
\caption{The carry predictor counts old payload reads, merge inputs, and
output rewrites. Cursor bitmap inspection is exported separately.}
\end{table}

\paragraph{Interpretation.}
Carry work, physical resolve/app work, seed work, and drain/synchronization
work explain their corresponding cycle counters with $R^2$ above 0.99 in the
covered population. Directory request count ($R^2=0.0011$) and the current
aggregate switch-work count ($R^2=0.0181$) do not transfer across mixed case
classes, so no explanatory claim is made for those two mechanisms yet.

\paragraph{Claim boundary.}
These correlations validate the simulator's internal cost model and expose
bottlenecks. They do not constitute cycle-for-cycle FPGA calibration, total
board energy, or proof that one scalar predictor explains every topology.
\end{document}
